\chapter{Measure-Theoretic Probability}

Having seen probability power modern applications in machine learning and Monte Carlo methods,
we now turn to the mathematical ground beneath them.

For thirty-five chapters we have explored probability across many settings:
dice and coins, densities and expectations, Markov chains and Brownian
motion, the law of large numbers and the central limit theorem. We have computed,
proved, simulated, and applied. One might reasonably ask: if standard probability
can be developed this far without measure theory, why introduce it at the end of the text?

The answer is that the elementary theory relies on an implicit assumption: that
every subset $A$ of the sample space can be assigned a probability $\P(A)$. For
finite and countable sample spaces, this assumption is mathematically sound. For
continuous sample spaces --- such as the unit interval $[0,1]$, the real line $\R$,
or the space of infinite coin-flip sequences --- it is false. As Giuseppe Vitali
proved in 1905, there exist subsets of $[0,1]$ to which no consistent notion of
length or probability can be assigned without producing mathematical contradictions:
for instance, assigning two different probabilities to the same set, or decomposing
a region into pieces that reassemble into something of a different total measure.

The continuum therefore requires a principled decision: either maintain separate,
ad hoc treatments for discrete and continuous probability, or specify precisely
\emph{which} subsets carry well-defined probabilities. Measure theory makes this
specification through three core concepts: $\sigma$-algebras, measures, and
measurable functions. Under this framework, discrete random variables and continuous
stochastic processes are instances of the same mathematical object, governed by
the same theorems. The discrete sum $\sum_x x\,\P(X=x)$ and the continuous integral
$\int x\,f(x)\,\dd x$ become two expressions of a single integral.

In this chapter we build those three ingredients and make ``almost surely'' precise;
Chapters~37 and~38 then construct the Lebesgue integral and conditional expectation on
top of them.

\section{The crack in the foundation}
\label{ch36:crack}

\begin{intuitionbox}
Probability is, at heart, the assignment of numbers to subsets of a sample space.
For dice, the subsets are finite and you can list them, so assigning numbers is no
trouble. But when the sample space is something like ``a point chosen uniformly at
random from $[0,1]$,'' the subsets form a wild, sprawling collection. The hidden
difficulty is that \emph{not every one of those subsets can carry a
sensible probability}. We have to find out which ones can, and restrict our
attention to those --- without losing anything we actually care about.
\end{intuitionbox}

Let us make the trouble concrete. Pick a point $U$ uniformly at random from the
interval $[0,1]$. ``Uniformly'' should mean: the probability of landing in a
sub-interval is its length. So $\P(U \in [0.2,0.5]) = 0.3$, and more generally
$\P(U \in [a,b]) = b-a$. This is the most natural continuous experiment there is,
the source of every density we ever drew. We would like a function $\P$ that eats
\emph{any} subset $A \subseteq [0,1]$ and returns its probability, agreeing with
length on intervals.

What properties must $\P$ have? Three properties appear indispensable, adapting
Kolmogorov's axioms from Chapter~3:

\begin{enumerate}
  \item \textbf{Non-negativity and normalization.} $\P(A) \ge 0$ for every $A$,
        and $\P([0,1]) = 1$.
  \item \textbf{Translation invariance.} Sliding a set sideways (wrapping around
        the circle, so $[0,1]$ becomes a loop) should not change its probability.
        A uniform point does not prefer the left half to the right.
  \item \textbf{Countable additivity.} If $A_1, A_2, \dots$ are disjoint, then
        $\P\!\left(\bigcup_n A_n\right) = \sum_n \P(A_n)$. The probability of a
        countable disjoint union is the sum of the probabilities. We \emph{relied}
        on this in every continuous-distribution argument in the book.
\end{enumerate}

These three look mild. The shocking fact, discovered by Giuseppe Vitali in 1905,
is that \emph{no function $\P$ defined on all subsets of $[0,1]$ can satisfy all
three at once.} Something has to give.

\begin{theorem}[There is no uniform probability on all subsets]
\label{ch36:vitali}
There is no function $\P$ assigning a number to \emph{every} subset of $[0,1]$
that is non-negative, normalized with $\P([0,1])=1$, translation-invariant (mod
$1$), and countably additive.
\end{theorem}

\begin{whybox}
The proof constructs a subset $V \subset [0,1)$ --- a \emph{Vitali set} --- and
shows that no numerical value can be consistently assigned to $\P(V)$ under the stated axioms.

Work on the circle $[0,1)$ with addition mod $1$. Declare two points $x,y$
\emph{equivalent} if their difference $x-y$ is rational. This partitions $[0,1)$ into
disjoint equivalence classes; each class is a copy of the rationals shifted by
some amount. Using the axiom of choice, pick exactly one representative from each
class and collect them into a set $V$.

Now list the rationals in $[0,1)$ as $q_0, q_1, q_2, \dots$ (they are countable),
and form the shifted copies $V_n = V + q_n \pmod 1$. Two facts hold:

\begin{itemize}
  \item The $V_n$ are \emph{pairwise disjoint}. If a point landed in both $V_n$
        and $V_m$, two representatives would differ by a rational, contradicting
        ``one representative per class.''
  \item The $V_n$ \emph{cover everything}: $\bigcup_n V_n = [0,1)$. Any point $x$
        is equivalent to the representative of its class, so $x$ differs from that
        representative by some rational $q_n$, putting $x$ in $V_n$.
\end{itemize}

So $[0,1)$ is the disjoint union of the countably many sets $V_n$, each a translation
of $V$. By translation invariance, $\P(V_n) = \P(V)$ for all $n$ --- call
this common value $c$. By countable additivity,
\[
  1 = \P([0,1)) = \sum_{n=0}^{\infty} \P(V_n) = \sum_{n=0}^{\infty} c .
\]
Now consider $\sum_n c$. If $c = 0$, the sum is $0 \ne 1$. If $c > 0$, the sum is
$+\infty \ne 1$. There is no real number $c$ for which a constant series of terms
equals $1$. This contradiction shows that no such probability function $\P$ can exist on all
subsets of $[0,1)$, and in particular $V$ cannot be assigned any probability consistent
with the axioms.
\end{whybox}

Sets like $V$ are called \emph{non-measurable}: if we keep countable additivity and
translation invariance on $[0,1]$, the subset $\{U \in V\}$ simply cannot be assigned a
probability.

\begin{connectionbox}
\textbf{The axiom of choice and Banach--Tarski.}
The axiom of choice, which we used to pick one representative from each of
uncountably many classes, was isolated by Ernst Zermelo in 1904. Vitali's
construction (1905) showed that the axiom has significant consequences: it
proves the existence of non-constructive sets that cannot be defined through
standard set operations on intervals. In 1924, Stefan Banach and Alfred Tarski
proved that using non-measurable subsets, a solid ball in three dimensions can
be decomposed into finitely many disjoint pieces and reassembled, via rigid rotations
and translations alone, into two solid balls of the same radius as the original.
Volume is not conserved because the individual pieces are non-measurable and have
no well-defined volume. Measure theory prevents such anomalies by restricting
the domain of measures to collections of measurable sets.
\end{connectionbox}

\begin{paradoxbox}
\textbf{``If a uniform random point lands somewhere, does it not have a probability of landing in $V$?''}
The outcome $U$ always takes some numerical value in $[0,1]$. However, the statement
``$U \in V$'' cannot be assigned a probability because $V$ is non-measurable.
Events in probability must be formed from intervals through countable unions,
intersections, and complements. The Vitali set requires the axiom of choice and lies
outside this collection. When the domain of $\P$ is restricted to a suitably defined
class of events, these contradictions are avoided, and probability theory operates rigorously.
\end{paradoxbox}

We cannot define $\P$ on all subsets of an uncountable sample space.
Instead, we define it on a carefully chosen collection of subsets --- large enough to
contain all intervals and their countable combinations, but restricted enough to exclude
non-measurable sets. That collection is a $\sigma$-algebra, the first of our three core objects.

\section{$\sigma$-algebras: the family of legitimate events}
\label{ch36:sigma-algebra}

\begin{intuitionbox}
A $\sigma$-algebra is a collection of subsets representing the events to which
probabilities can be consistently assigned. If an event $A$ is measurable, its
complement $A^c$ (``not $A$'') must also be measurable. If $A$ and $B$ are events,
so are their union $A \cup B$ (``$A$ or $B$'') and intersection $A \cap B$ (``$A$ and $B$'')
--- and this closure extends to countably many events simultaneously.
\end{intuitionbox}

Let $\Omega$ be a sample space --- the set of all outcomes. We want to single out
a family $\mathcal{F}$ of subsets of $\Omega$, to be called the \emph{events},
such that we can manipulate events the way probability requires.

\begin{definition}[$\sigma$-algebra]
\label{ch36:sigma-def}
A collection $\mathcal{F}$ of subsets of $\Omega$ is a \emph{$\sigma$-algebra} (or
\emph{$\sigma$-field}) if:
\begin{enumerate}
  \item $\Omega \in \mathcal{F}$ \quad(the certain event is allowed);
  \item if $A \in \mathcal{F}$ then $A^c = \Omega \setminus A \in \mathcal{F}$
        \quad(closed under complement);
  \item if $A_1, A_2, \dots \in \mathcal{F}$ then
        $\bigcup_{n=1}^{\infty} A_n \in \mathcal{F}$
        \quad(closed under \emph{countable} unions).
\end{enumerate}
The pair $(\Omega, \mathcal{F})$ is called a \emph{measurable space}, and the
members of $\mathcal{F}$ are the \emph{measurable sets} or \emph{events}.
\end{definition}

Three short axioms, and they automatically deliver everything else you would want.

\begin{proposition}[What a $\sigma$-algebra contains for free]
\label{ch36:closure}
If $\mathcal{F}$ is a $\sigma$-algebra then:
(i) $\emptyset \in \mathcal{F}$;
(ii) $\mathcal{F}$ is closed under \emph{countable intersections};
(iii) $\mathcal{F}$ is closed under set differences $A \setminus B$.
\end{proposition}

\begin{whybox}
Each follows by combining complement with union, the way De~Morgan's laws let you
trade ``and'' for ``or.''
(i) $\Omega \in \mathcal{F}$, so $\emptyset = \Omega^c \in \mathcal{F}$.
(ii) Given $A_n \in \mathcal{F}$, each complement $A_n^c \in \mathcal{F}$, their
union $\bigcup_n A_n^c \in \mathcal{F}$, and its complement is, by De~Morgan,
$\bigcap_n A_n \in \mathcal{F}$.
(iii) $A \setminus B = A \cap B^c$, an intersection of two members.
Notice that while the definition specifies only complements and countable
unions, De~Morgan's laws guarantee closure under countable intersections and set
differences as well. The axioms provide a minimal basis that ensures closure
under all standard countable set operations.
\end{whybox}

\begin{focusbox}
Why \emph{countable} unions, rather than finite or arbitrary?
Closure under \emph{finite} unions alone is insufficient for continuous probability,
which relies on limiting operations. Events such as $\set{X \le x}$ are frequently
expressed as $\bigcap_n \set{X < x + 1/n}$, a countable intersection. Conversely,
closure under \emph{arbitrary} unions is too strong: every singleton $\set{x}$ is a Borel
set, and an arbitrary union of singletons forms an arbitrary subset of $\R$ --- including
the non-measurable Vitali set. Allowing arbitrary unions would make every subset measurable,
reintroducing the contradiction of Theorem~\ref{ch36:vitali}. Countable operations provide
the exact structure needed for limits in analysis while preserving mathematical consistency.
The prefix ``$\sigma$-'' denotes closure under countable operations.
\end{focusbox}

\subsection{Examples, from trivial to indispensable}

\begin{example}[The two extremes]
\label{ch36:extremes}
The smallest $\sigma$-algebra on any $\Omega$ is the \emph{trivial} one
$\mathcal{F} = \set{\emptyset, \Omega}$.
The largest is the \emph{power set} $\mathcal{F} = 2^{\Omega}$, containing all subsets.
For finite or countable $\Omega$, the power set is the standard choice and non-measurable
sets do not arise. This is why elementary discrete probability can be developed without
measure theory; the need to restrict the collection of events arises specifically on
uncountable sample spaces.
\end{example}

\begin{example}[Generating from a question]
\label{ch36:generated}
Often you have a few sets you definitely want to be events --- say a single
event $A$ --- and you want the smallest $\sigma$-algebra containing them. For one
set $A$ it is $\set{\emptyset, A, A^c, \Omega}$: the four answers to ``what is the
status of $A$?'' In general, given any collection $\mathcal{C}$ of subsets, the
\emph{$\sigma$-algebra generated by $\mathcal{C}$}, written $\sigma(\mathcal{C})$,
is the intersection of all $\sigma$-algebras containing $\mathcal{C}$. (An
intersection of $\sigma$-algebras is again a $\sigma$-algebra --- check the three
axioms --- so this is well defined and is the smallest one that works.)

A helpful mental model for a sub-$\sigma$-algebra $\mathcal{G} \subseteq \mathcal{F}$ is the \emph{resolving power of an optical microscope or digital sensor} (Figure~\ref{ch36:fig-partition-tree}). Think of the sample space $\Omega$ as a canvas. When your lens has zero resolving power ($\mathcal{F}_0 = \{\emptyset, \Omega\}$), the entire canvas is a single blurred pixel: you can only tell whether something happened or didn't. When you make an observation---say, revealing the first coin toss $X_1$---the sensor's resolution doubles to $\mathcal{F}_1 = \sigma(X_1)$, cleaving the canvas into two distinct pixel blocks $A_H$ and $A_T$. You can now distinguish whether an outcome fell in $A_H$ or $A_T$, and the measurable events are all unions of these blocks. If a subsequent coin toss $X_2$ is observed, each block splits again into four finer pixels in $\mathcal{F}_2 = \sigma(X_1, X_2)$.
A set is measurable with respect to a sub-$\sigma$-algebra if and only if its boundary aligns with the grain of the sensor's pixels. If you try to ask a question whose truth depends on finer structure than the current resolution---such as asking at stage $\mathcal{F}_1$ whether the \emph{second} toss landed heads---the question cuts straight across the indivisible pixels without containing them whole. The sensor simply cannot resolve it.
\end{example}

\begin{figure}[H]
\centering
\begin{tikzpicture}[>=Stealth, font=\small]
% Lane 1: F0 (Trivial resolution)
\begin{scope}[yshift=4.0cm]
  \node[anchor=east, font=\footnotesize\bfseries] at (-0.4, 0.5) {$\mathcal{F}_0 = \{\emptyset, \Omega\}$};
  \draw[thick, draw=sourceblue!80!black, fill=sourceblue!12, rounded corners=2pt] (0,0) rectangle (7,1);
  \node[font=\footnotesize] at (3.5, 0.5) {$\Omega$ (undivided sample space: 1 pixel block)};
  \node[anchor=west, font=\scriptsize, sourceblue!80!black] at (7.3, 0.5) {Macro-state: zero information};
\end{scope}

% Downward arrow 1
\draw[->, thick, gray!70] (3.5, 3.8) -- (3.5, 3.0) node[midway, right, font=\tiny, gray!90] {observe flip 1};

% Lane 2: F1 (Coarse resolution)
\begin{scope}[yshift=1.8cm]
  \node[anchor=east, font=\footnotesize\bfseries] at (-0.4, 0.5) {$\mathcal{F}_1 = \sigma(X_1)$};
  \draw[thick, draw=targetgreen!80!black, fill=targetgreen!15, rounded corners=2pt] (0,0) rectangle (3.4,1);
  \draw[thick, draw=quotientorange!80!black, fill=quotientorange!15, rounded corners=2pt] (3.6,0) rectangle (7,1);
  \node[font=\footnotesize] at (1.7, 0.5) {Block $A_H$: $\{H**\}$};
  \node[font=\footnotesize] at (5.3, 0.5) {Block $A_T$: $\{T**\}$};
  \node[anchor=west, font=\scriptsize, targetgreen!80!black] at (7.3, 0.5) {2 pixel blocks ($2^2=4$ events)};
\end{scope}

% Downward arrow 2
\draw[->, thick, gray!70] (3.5, 1.6) -- (3.5, 0.8) node[midway, right, font=\tiny, gray!90] {observe flip 2};

% Lane 3: F2 (Finer resolution)
\begin{scope}[yshift=-0.4cm]
  \node[anchor=east, font=\footnotesize\bfseries] at (-0.4, 0.5) {$\mathcal{F}_2 = \sigma(X_1, X_2)$};
  \draw[thick, draw=structurepurple!80!black, fill=structurepurple!12, rounded corners=1pt] (0,0) rectangle (1.6,1);
  \draw[thick, draw=structurepurple!80!black, fill=structurepurple!18, rounded corners=1pt] (1.8,0) rectangle (3.4,1);
  \draw[thick, draw=quotientorange!80!black, fill=quotientorange!12, rounded corners=1pt] (3.6,0) rectangle (5.2,1);
  \draw[thick, draw=quotientorange!80!black, fill=quotientorange!18, rounded corners=1pt] (5.4,0) rectangle (7,1);

  \node[font=\scriptsize] at (0.8, 0.5) {$\{HH*\}$};
  \node[font=\scriptsize] at (2.6, 0.5) {$\{HT*\}$};
  \node[font=\scriptsize] at (4.4, 0.5) {$\{TH*\}$};
  \node[font=\scriptsize] at (6.2, 0.5) {$\{TT*\}$};
  \node[anchor=west, font=\scriptsize, structurepurple!80!black] at (7.3, 0.5) {4 pixel blocks ($2^4=16$ events)};
\end{scope}

% Bottom summary callout
\node[font=\footnotesize, text width=10.5cm, align=center, draw=gray!30, fill=gray!5, rounded corners=3pt, inner sep=6pt] at (3.5, -1.3)
  {Filtration $\mathcal{F}_0 \subset \mathcal{F}_1 \subset \mathcal{F}_2$: a sub-$\sigma$-algebra is the \emph{resolving power} of an optical sensor.\\
   An event is measurable with respect to $\mathcal{F}_t$ if and only if it is a union of stage-$t$ pixel blocks.};
\end{tikzpicture}
\caption{The $\sigma$-algebra as the resolving power of a microscope. At the coarsest level $\mathcal{F}_0$, the sensor sees only the entire space $\Omega$ (no individual features can be discerned). Observing the first coin toss $X_1$ doubles the resolution to $\mathcal{F}_1$, cleaving the space into two atomic blocks $A_H$ and $A_T$. Observing the second toss $X_2$ refines each block into two sub-blocks in $\mathcal{F}_2$. A question like ``did the second toss land heads?'' cannot be answered at stage $\mathcal{F}_1$ because the event $\{HH*, TH*\}$ cuts across the grain of individual pixel blocks without containing them whole.}
\label{ch36:fig-partition-tree}
\end{figure}

The single most important generated $\sigma$-algebra is the one that tames the
real line.

\begin{definition}[Borel $\sigma$-algebra]
\label{ch36:borel}
The \emph{Borel $\sigma$-algebra} $\mathcal{B}(\R)$ is the $\sigma$-algebra on
$\R$ generated by the open intervals (equivalently, by the half-lines
$(-\infty, x]$). Its members are the \emph{Borel sets}.
\end{definition}

\begin{happeningbox}
The Borel sets encompass every set obtainable from intervals through countable
unions, countable intersections, and complements. This includes all open and closed
sets, individual points, the rationals, the irrationals, and the Cantor set.
Constructing a non-Borel subset of $\R$ requires non-constructive principles such
as the axiom of choice. Restricting events to the Borel $\sigma$-algebra $\mathcal{B}(\R)$
accommodates all sets encountered in analysis while excluding non-measurable constructions.
Every density, distribution function, and interval probability $\P(a \le X \le b)$
discussed in earlier chapters represents a statement about Borel sets.
\end{happeningbox}

\begin{connectionbox}
There is a direct analogy with linear algebra (see the companion linear-algebra
volume). A $\sigma$-algebra generated by a collection $\mathcal{C}$ is the
\emph{closure} of $\mathcal{C}$ under countable set operations, just as the
\emph{span} of a set of vectors is the closure under linear combinations, or the
\emph{subgroup generated} by a subset is the closure under the group operation.
``Generated by'' denotes the smallest structure of a given type containing the
generators, obtained by intersecting all such containing structures.
\end{connectionbox}

\section{Measures and probability measures}
\label{ch36:measures}

We now have the sample space $\Omega$ and the $\sigma$-algebra of allowed events $\mathcal{F}$.
The final ingredient is the assignment of measure. A \emph{measure} formalizes the
notion of size --- length, area, volume, mass, or probability --- satisfying countable additivity.

\begin{definition}[Measure and probability measure]
\label{ch36:measure-def}
Let $(\Omega, \mathcal{F})$ be a measurable space. A \emph{measure} is a function
$\mu : \mathcal{F} \to [0, \infty]$ such that
\begin{enumerate}
  \item $\mu(\emptyset) = 0$;
  \item \textbf{Countable additivity (}$\sigma$\textbf{-additivity):} for disjoint
        $A_1, A_2, \dots \in \mathcal{F}$,
        \[
          \mu\!\left(\bigcup_{n=1}^{\infty} A_n\right)
          = \sum_{n=1}^{\infty} \mu(A_n).
        \]
\end{enumerate}
If in addition $\mu(\Omega) = 1$, we call $\mu$ a \emph{probability measure} and
usually write it $\P$. The triple $(\Omega, \mathcal{F}, \P)$ is a
\emph{probability space} --- the complete description of a random experiment.
\end{definition}

These conditions formalize Kolmogorov's axioms: non-negativity
($\mu \ge 0$ is built into the codomain), normalization ($\P(\Omega)=1$), and
countable additivity. The measure-theoretic framework formulates these axioms
rigorously for arbitrary sample spaces.

\begin{historybox}
\textbf{Kolmogorov, 1933: an axiomatic home for uncertainty.}
In 1933, Andrey Kolmogorov published \emph{Grundbegriffe der Wahrscheinlichkeitsrechnung}
(``Foundations of the Theory of Probability''), a slim monograph of about sixty pages that
did for probability what Euclid did for geometry. Before it, probability was a collection
of clever techniques without an agreed foundation; mathematicians argued over what a
random variable even was, and paradoxes on the continuum lingered unresolved.

Kolmogorov's insight was that the measure theory developed by \'Emile Borel and Henri
Lebesgue for length and area was \emph{the same theory} needed for probability: normalize
total mass to $1$, restrict events to a $\sigma$-algebra, and rename ``measurable set'' as
``event.'' At a stroke, every theorem about Lebesgue integration became a theorem about
expectation. Richard von Mises and others initially objected that an axiomatic system was
detached from empirical frequencies, but Kolmogorov's framework quickly became the standard
mathematical foundation for modern probability and stochastic processes.
\end{historybox}

\subsection{Three measures you already know}

\begin{example}[Counting measure: discrete probability]
\label{ch36:counting}
Take $\Omega$ countable, $\mathcal{F} = 2^{\Omega}$, and a function
$p : \Omega \to [0,1]$ with $\sum_{\omega} p(\omega) = 1$ --- a pmf. Define
$\P(A) = \sum_{\omega \in A} p(\omega)$. Countable additivity is just the fact
that you can split a sum over a disjoint union into separate sums. This is
\emph{all of discrete probability}, repackaged. A fair die is
$\Omega = \set{1,\dots,6}$, $p(\omega) = 1/6$.
\end{example}

\begin{example}[Lebesgue measure: length, area, the continuum]
\label{ch36:lebesgue}
On $(\R, \mathcal{B}(\R))$ there is a unique measure $\lambda$ with
$\lambda([a,b]) = b - a$ for every interval: \emph{Lebesgue measure}, the rigorous
notion of length. It assigns $\lambda(\set{x}) = 0$ to every point (a point has no
length), and --- a fact that surprises everyone --- $\lambda(\Q) = 0$: the
rationals, dense though they are, have total length zero, because they are
countable and you can cover the $n$th rational by an interval of length
$\eps/2^n$, for total length $\eps$, as small as you like. Restricting $\lambda$
to $[0,1]$ gives the uniform probability measure we tried and failed to define on
\emph{all} subsets in Section~\ref{ch36:crack} --- it works perfectly once we
confine it to the Borel sets.
\end{example}

\begin{example}[Dirac measure: a point mass]
\label{ch36:dirac}
Fix a point $a$. The \emph{Dirac measure} $\delta_a$ assigns $\delta_a(A) = 1$ if
$a \in A$ and $0$ otherwise --- all the mass piled at a single point. As a
probability measure it is a degenerate random variable, ``equals $a$ with
certainty.'' Mixing Dirac measures with Lebesgue measure is how a single framework
handles random variables that are part-discrete, part-continuous
(Section~\ref{ch36:unify}).
\end{example}

\subsection{Properties that fall out for free}

The same countable additivity that defines a measure forces the everyday
identities you have used all along. We prove a representative few; they are
templates for the rest.

\begin{proposition}[Basic measure properties]
\label{ch36:props}
For a probability measure $\P$ and events $A, B$:
\begin{enumerate}
  \item \textbf{Complement:} $\P(A^c) = 1 - \P(A)$.
  \item \textbf{Monotonicity:} if $A \subseteq B$ then $\P(A) \le \P(B)$.
  \item \textbf{Inclusion--exclusion:} $\P(A \cup B) = \P(A) + \P(B) - \P(A\cap B)$.
  \item \textbf{Continuity from below:} if $A_1 \subseteq A_2 \subseteq \cdots$
        then $\P\!\left(\bigcup_n A_n\right) = \lim_{n\to\infty} \P(A_n)$.
\end{enumerate}
\end{proposition}

\begin{whybox}
For (1), $\Omega = A \cup A^c$ disjointly, so additivity gives
$1 = \P(\Omega) = \P(A) + \P(A^c)$.
For (2), write $B = A \cup (B \setminus A)$ disjointly; then
$\P(B) = \P(A) + \P(B\setminus A) \ge \P(A)$ because the extra term is $\ge 0$.
For (3), split $A \cup B$ into the three disjoint pieces
$A\setminus B$, $A\cap B$, $B\setminus A$ and add.

Property (4), \emph{continuity}, is an essential consequence of \emph{countable}
(rather than merely finite) additivity, allowing probability to commute
with limits along monotone sequences of events.
Set $B_1 = A_1$ and $B_n = A_n \setminus A_{n-1}$ for $n \ge 2$. The $B_n$ are
disjoint, $\bigcup_{k\le n} B_k = A_n$, and $\bigcup_k B_k = \bigcup_n A_n$. By
countable additivity,
\[
  \P\!\left(\bigcup_n A_n\right) = \sum_{k=1}^{\infty} \P(B_k)
  = \lim_{n\to\infty} \sum_{k=1}^{n} \P(B_k)
  = \lim_{n\to\infty} \P(A_n).
\]
The limit corresponds directly to the infinite series, demonstrating that
continuity of a measure along nested sequences is equivalent to countable additivity.
\end{whybox}

\begin{figure}[H]
\centering
\begin{tikzpicture}
\begin{axis}[
  width=12cm, height=6cm,
  xlabel={$n$ (number of nested events)},
  ylabel={$\P(A_n)$},
  xmin=0, xmax=12, ymin=0, ymax=1.14,
  domain=1:12, samples=12,
  title={Continuity from below: $\P(A_n)\to\P(\bigcup_n A_n)$},
  title style={font=\small},
  axis lines=left, grid=both, grid style={gray!18},
  ymajorgrids=true,
  ]
\addplot[ycomb, mark=*, mark size=1.6pt, color=distblue, thick]
  {1 - 0.65*exp(-0.45*x)};
\addplot[dashed, densityred, thick, domain=0:12] {1};
\node[densityred, anchor=south east] at (axis cs:11.8,1.01) {limit $\P\big(\bigcup_n A_n\big)=1$};
\end{axis}
\end{tikzpicture}
\caption{As the nested events $A_n$ expand toward their union, their probabilities
converge to the probability of the union. Countable additivity guarantees that this
limit is attained; finite additivity alone would not suffice.}
\label{ch36:fig-continuity}
\end{figure}

Figure~\ref{ch36:fig-continuity} plots this convergence for an expanding sequence of events;
monotone limit arguments throughout probability rely on it.

\section{Random variables, done right}
\label{ch36:random-variables}

In Chapter~6 we introduced a random variable as a function from the sample space to
the real numbers. To make this definition rigorous on general probability spaces, one
further condition is required: the function must be \emph{measurable}.

\begin{intuitionbox}
A random variable $X$ assigns a real number to each outcome $\omega \in \Omega$.
To evaluate probabilities such as $\P(X \in B)$, the preimage
$X^{-1}(B) = \set{\omega : X(\omega) \in B}$ must belong to the domain of $\P$,
which is $\mathcal{F}$. Measurability guarantees that preimages of Borel sets are
valid events.
\end{intuitionbox}

\begin{definition}[Measurable function / random variable]
\label{ch36:rv-def}
Let $(\Omega, \mathcal{F})$ be a measurable space. A function $X : \Omega \to \R$
is \emph{measurable} (a \emph{random variable}) if for every Borel set
$B \in \mathcal{B}(\R)$,
\[
  X^{-1}(B) = \set{\omega \in \Omega : X(\omega) \in B} \in \mathcal{F}.
\]
Equivalently --- and this is the version you check in practice --- it suffices that
$\set{\omega : X(\omega) \le x} \in \mathcal{F}$ for every $x \in \R$.
\end{definition}

\begin{whybox}
Why is checking the half-lines $\set{X \le x}$ enough to guarantee \emph{all}
Borel sets pull back into $\mathcal{F}$? Because preimages respect set operations
perfectly: $X^{-1}(B^c) = (X^{-1}(B))^c$ and
$X^{-1}\!\left(\bigcup_n B_n\right) = \bigcup_n X^{-1}(B_n)$. So the collection of
sets $B$ whose preimage lands in $\mathcal{F}$ is \emph{itself} a $\sigma$-algebra.
If that collection contains the half-lines $(-\infty, x]$, it must contain the
entire $\sigma$-algebra they generate --- which is $\mathcal{B}(\R)$. This relies
on the principle that verifying measurability on a generating class guarantees
measurability on the entire generated $\sigma$-algebra.
\end{whybox}

\begin{happeningbox}
Measurability formalizes the concept of \emph{information}. A $\sigma$-algebra
$\mathcal{F}$ represents the collection of events to which probabilities can be assigned;
a function is measurable when every Borel question regarding its values pulls back to an
event in $\mathcal{F}$. This interpretation is especially useful with \emph{sub-}$\sigma$-algebras:
a sub-$\sigma$-algebra $\mathcal{G} \subseteq \mathcal{F}$ models \emph{partial} information,
representing the events discernible from a subset of observations. A random variable is
$\mathcal{G}$-measurable precisely when its value can be determined from that partial
information. This provides the mathematical basis for conditioning, filtrations in martingale
theory (Chapter~28), and information structures in stochastic processes.
\end{happeningbox}

\begin{example}[Standard functions are measurable]
\label{ch36:meas-examples}
Virtually all functions encountered in analysis and applications are measurable.
Constants are measurable. Indicators $\ind_A$ are measurable if and only if $A \in \mathcal{F}$.
Pointwise sums, products, maxima, minima, and limits of sequences of measurable functions
are measurable. Continuous functions $\R \to \R$, monotone functions, and piecewise-continuous
functions are Borel measurable. Constructing a non-measurable function requires non-constructive
tools such as the axiom of choice; the condition serves to ensure analytical consistency.
\end{example}

\subsection{The distribution: pushing the measure forward}

A random variable carries the probability measure on $\Omega$ over to a
probability measure on $\R$, and that induced measure is what earlier chapters
computed with.

\begin{definition}[Distribution / law]
\label{ch36:law}
The \emph{distribution} (or \emph{law}) of a random variable $X$ is the
probability measure $\P_X$ on $(\R, \mathcal{B}(\R))$ defined by
\[
  \P_X(B) = \P\big(X^{-1}(B)\big) = \P(X \in B),
  \qquad B \in \mathcal{B}(\R).
\]
\end{definition}

This is the \emph{pushforward} of $\P$ by $X$. It is a probability
measure on $\R$ --- measurability ensures that $X^{-1}(B) \in \mathcal{F}$
for every Borel set $B$, and countable additivity transfers because preimages
preserve disjoint unions.

\begin{aestheticbox}
Two random variables can be defined on completely different probability spaces ---
one from coin tosses, another from an interval with Lebesgue measure --- yet induce
the \emph{same} distribution $\P_X$ on $\R$. In earlier chapters, stating
``let $X \sim \Normal(0,1)$'' required no specification of the underlying sample space
$\Omega$. The probabilistic properties of $X$ depend entirely on the pushforward
measure $\P_X$ on $(\R, \mathcal{B}(\R))$ rather than the particular construction of $\Omega$.
\end{aestheticbox}

\begin{connectionbox}
The cumulative distribution function $F_X(x) = \P_X((-\infty, x]) = \P(X \le x)$ is
the pushforward measure evaluated on half-lines. Because the collection of half-lines
generates $\mathcal{B}(\R)$, the CDF uniquely determines the entire probability measure
$\P_X$. This provides the rigorous justification for using the CDF as a complete
characterization of a univariate distribution.
\end{connectionbox}

\section{``Almost surely'': events of probability one}
\label{ch36:almost-surely}

You have seen the phrase ``almost surely'' (a.s.) and ``with probability $1$''
throughout the convergence chapters. Measure theory gives it an exact and
surprisingly subtle meaning.

\begin{definition}[Null sets and almost-sure events]
\label{ch36:null}
An event $N$ is a \emph{null set} if $\P(N) = 0$. A statement holds
\emph{almost surely} (or \emph{almost everywhere}, in pure measure theory) if the
set of outcomes where it \emph{fails} is contained in a null set --- i.e.\ it holds
except on a set of probability $0$.
\end{definition}

The subtlety is that \emph{``probability zero'' is not the same as
``impossible.''} On $\Omega = [0,1]$ with uniform measure, the event
$\set{U = 0.5}$ is a single point, so $\P(U = 0.5) = 0$ --- yet it is a perfectly
possible outcome; the point $0.5$ is right there in the sample space. Indeed
\emph{every} particular value the random point takes had probability $0$ before
the experiment, and yet some value occurs. Probability $0$ means
``negligible in aggregate,'' not ``cannot happen.''

\begin{paradoxbox}
\textbf{The dartboard paradox.} Throw a dart at the interval $[0,1]$; it lands at a
uniformly random point $U$. The probability it hits any specific point is $0$. So
the probability it hits \emph{some} point should be ``a sum of zeros''
$=0$ --- yet the dart certainly lands somewhere. Where did the probability come
from? \emph{Resolution:} the union $[0,1] = \bigcup_{x \in [0,1]} \set{x}$ is an
\emph{uncountable} union, and additivity --- even countable additivity --- says
\emph{nothing} about uncountable sums of zeros. Countable additivity gives
$\P(\Q \cap [0,1]) = \sum_{q} 0 = 0$, correctly, because the rationals are
countable. But you may not add up the uncountably many points of $[0,1]$ that way.
The dartboard paradox is, at root, a misuse of additivity across an uncountable
index set --- the very abuse the $\sigma$ in $\sigma$-algebra was designed to
forbid. The continuum holds positive probability precisely because it cannot be
disassembled into countably many zero-probability atoms.
\end{paradoxbox}

\begin{focusbox}
Why ``almost surely'' is the natural mode of truth in continuous probability:
because changing a random variable on a null set changes \emph{nothing} you can
measure. If $X = Y$ except on a set of probability $0$, then
$\P(X \in B) = \P(Y \in B)$ for every $B$, so $X$ and $Y$ have identical
distributions, identical expectations, identical everything observable. Measure theory therefore identifies functions that agree almost surely.
Theorems in modern probability are naturally stated in terms of almost-sure
convergence: requiring convergence on every outcome without exception is
unnecessarily restrictive when discrepancies on null sets do not affect
probabilities or expectations.
\end{focusbox}

\begin{example}[Strong law, reread]
\label{ch36:sln-reread}
The strong law of large numbers (Chapter~20) says
$\bar{X}_n \asto \mu$. Decoded: the set of infinite outcome sequences
$\omega = (\omega_1, \omega_2, \dots)$ for which the running average
$\bar{X}_n(\omega)$ \emph{fails} to converge to $\mu$ is a null set. Sequences
that misbehave certainly \emph{exist} --- ``always heads forever'' is one --- but
together they form a set of probability $0$. The strong law is, word for word, an
almost-sure statement, and it can only even be \emph{phrased} on a probability
space rich enough to host infinite sequences, namely
$\Omega = \set{H,T\,}^{\N}$ with a product $\sigma$-algebra. Measure theory is
what lets us build that space and put a coherent probability on it.
\end{example}

\section{One framework, two worlds: unifying discrete and continuous}
\label{ch36:unify}

Measure theory provides a unified framework encompassing both discrete and continuous
probability.

In Chapter~8, expectation was defined separately as a sum for discrete variables,
$\E[X] = \sum_x x\,\P(X=x)$, and as an integral for continuous variables,
$\E[X] = \int x\,f(x)\,\dd x$. We proved key properties (such as linearity and LOTUS)
twice. The Lebesgue integral introduced in the next chapter unifies these definitions:
both are instances of the single abstract integral
\[
  \E[X] = \int_{\Omega} X \,\dd\P,
\]
the integral of the function $X$ against the probability measure $\P$. When $\P$
is a discrete counting measure, this integral reduces to a sum; when $\P$ has a
density with respect to Lebesgue measure, it reduces to the standard integral.

The bridge that names the density rigorously is the Radon--Nikodym theorem, and we
can state its meaning now even before the next chapter proves it.

\begin{focusbox}
A distribution $\P_X$ on $\R$ has a \emph{density} $f$ exactly when it is
``compatible with length'' --- when every set of \emph{zero length} also has
\emph{zero probability} (this compatibility is called \emph{absolute continuity}
with respect to Lebesgue measure). In that case $f$ is the local exchange rate
between probability and length: $\P_X(B) = \int_B f\,\dd\lambda$. A distribution is
\emph{discrete} when, by contrast, all its mass sits on a countable set of atoms
(a sum of Dirac measures). And nothing forces a distribution to be one or the
other.
\end{focusbox}

\begin{example}[A variable that is neither discrete nor continuous]
\label{ch36:mixed}
Let $T$ be the time you wait at a crosswalk. With probability $1/2$ the light is
already green and $T = 0$; otherwise $T$ is uniform on $(0,1)$ minutes. Then
\[
  \P_T = \tfrac{1}{2}\,\delta_0 + \tfrac{1}{2}\,\lambda|_{(0,1)},
\]
half a point mass at $0$ plus half a stretch of uniform density. This $T$ has no
pmf (it spreads mass over an interval) and no density (it puts a positive lump on
the single point $0$, where any density would integrate to $0$). Elementary
probability typically treats discrete and continuous variables separately, leaving
mixed distributions as awkward hybrids. In measure theory, mixed variables require
no special treatment: a measure can decompose into discrete and continuous components,
and $\E[T] = \int_{\Omega} T\,\dd\P$ applies directly to both.
\end{example}

\begin{figure}[H]
\centering
\begin{tikzpicture}
\begin{axis}[
  width=12cm, height=6cm,
  xlabel={$t$ (minutes waited)},
  ylabel={probability content},
  xmin=-0.4, xmax=1.4, ymin=0, ymax=0.62,
  axis lines=left, grid=both, grid style={gray!16},
  title={A mixed distribution: an atom plus a slab of density},
  title style={font=\small},
  ytick={0,0.5}, xtick={0,0.5,1},
]
% density part: uniform 1/2 on (0,1)
\addplot[draw=distblue, thick, fill=distblue!20, domain=0:1, samples=2]
  coordinates {(0,0.5) (1,0.5)} \closedcycle;
% atom at 0
\addplot[ycomb, mark=*, mark size=2.6pt, color=densityred, very thick]
  coordinates {(0,0.5)};
\node[densityred, anchor=west] at (axis cs:0.04,0.55) {atom $1/2\,\delta_0$};
\node[distblue, anchor=south] at (axis cs:0.5,0.25)
  {density $1/2$ on $(0,1)$};
\end{axis}
\end{tikzpicture}
\caption{The waiting time $T$ of Example~\ref{ch36:mixed}. Half the probability
is a point mass (the red spike at $t=0$: ``light already green''), half is spread
as a uniform density over $(0,1)$ (the blue slab). Neither a pmf nor a pdf alone
describes this, but the measure $\P_T$ characterizes it completely.}
\label{ch36:fig-mixed}
\end{figure}

The same mixed law looks even cleaner through its CDF, where the atom and the
spread leave unmistakable signatures.

\begin{figure}[H]
\centering
\begin{tikzpicture}
\begin{axis}[
  width=12cm, height=6cm,
  xlabel={$t$},
  ylabel={$F_T(t)=\P(T\le t)$},
  xmin=-0.6, xmax=1.6, ymin=-0.05, ymax=1.1,
  axis lines=left, grid=both, grid style={gray!16},
  title={CDF of the mixed waiting time $T$},
  title style={font=\small},
  ytick={0,0.5,1}, xtick={0,0.5,1},
]
% F=0 for t<0
\addplot[distblue, very thick, domain=-0.6:0, samples=2] {0};
% jump segment at 0 (dashed vertical to show the discontinuity)
\addplot[densityred, dashed, thick] coordinates {(0,0) (0,0.5)};
% F = 1/2 + t/2 on (0,1)
\addplot[distblue, very thick, domain=0:1, samples=2] {0.5 + 0.5*x};
% F=1 for t>=1
\addplot[distblue, very thick, domain=1:1.6, samples=2] {1};
% open/closed dots at the jump
\filldraw[distblue] (axis cs:0,0.5) circle (2.2pt);
\draw[distblue, thick, fill=white] (axis cs:0,0) circle (2.2pt);
\node[densityred, anchor=west] at (axis cs:0.06,0.27) {jump of $1/2$ (the atom)};
\node[distblue, anchor=north west] at (axis cs:0.08,0.92) {slope $1/2$ (the density)};
\end{axis}
\end{tikzpicture}
\caption{The cumulative distribution function of $T$. The jump of height
$\tfrac12$ at $t=0$ is the point mass $\tfrac12\,\delta_0$; the steady rise of
slope $\tfrac12$ across $(0,1)$ is the uniform density. A pure pmf would give a
staircase of jumps; a pure density would give a smooth curve with no jumps. A
mixed law shows both at once, and the CDF --- which always exists --- is the one
description that copes with every distribution. (Exercise~\ref{ch36:ex-mixed}
asks you to derive this curve.)}
\label{ch36:fig-mixed-cdf}
\end{figure}

\begin{applicationbox}
\textbf{Why machine learning quietly runs on measures.} A modern generative model
defines a probability distribution over a high-dimensional space (images, text
embeddings, sensor readings). Loss functions like the KL divergence
$\KL(P \,\|\, Q) = \int \log\frac{\dd P}{\dd Q}\,\dd P$ are written with exactly the
Radon--Nikodym derivative $\dd P/\dd Q$ of this chapter; it only makes sense when
$P$ is absolutely continuous with respect to $Q$ (else the divergence is
$+\infty$, the model's way of saying ``$Q$ rules out something $P$ insists on'').
Diffusion models, normalizing flows, and variational inference all manipulate
densities-relative-to-a-reference --- pushforwards of a base measure through a
learned map, scored by change-of-variables. Importance sampling reweights one
measure to estimate expectations under another via the same derivative. The
engineer rarely says ``measure,'' but every change-of-distribution step is a
measure-theoretic operation, and the framework of this chapter is what guarantees
the manipulations are legal. The same goes for signal processing (see the
companion volume): a power spectrum is a measure on frequency, and white noise is
rigorously a random \emph{measure}, not a function.
\end{applicationbox}

\section{Putting it together: the anatomy of a probability space}
\label{ch36:anatomy}

Let us assemble the three objects into the single picture they were always meant to
form. Every random experiment in this book --- and in all of probability --- is a
triple
\[
  (\Omega, \mathcal{F}, \P),
\]
the \emph{probability space}, on which random variables live as measurable
functions:

\begin{itemize}
  \item $\Omega$ --- the \emph{sample space}, the set of all outcomes $\omega$.
        Often left unspecified, because only the induced laws matter.
  \item $\mathcal{F}$ --- the \emph{$\sigma$-algebra} of events, the collection of
        subsets closed under complement and countable union to which probabilities
        can be consistently assigned.
  \item $\P$ --- the \emph{probability measure}, assigning each event a number in
        $[0,1]$, countably additive, with $\P(\Omega) = 1$.
\end{itemize}

A random variable $X : \Omega \to \R$ is a measurable function; its law
$\P_X = \P \circ X^{-1}$ is a probability measure on $\R$; its CDF, density or pmf,
expectation, and variance are all read off from $\P_X$.

\begin{figure}[H]
\centering
\begin{tikzpicture}[
  every node/.style={font=\small},
  blob/.style={draw, thick, ellipse, minimum width=3.3cm, minimum height=2.1cm},
  realline/.style={draw, thick},
  arr/.style={-{Stealth[length=3mm]}, thick},
]
% Omega blob
\node[blob, fill=sourceblue!10, draw=sourceblue] (omega) at (0,0) {};
\node[sourceblue] at (0,1.35) {$(\Omega,\mathcal{F},\P)$};
\node[sourceblue] at (0,-1.35) {sample space};
\filldraw[sourceblue] (-0.6,0.1) circle (1.2pt) node[above, font=\tiny]{$\omega$};
\filldraw[sourceblue] (0.4,-0.3) circle (1.2pt);
\filldraw[sourceblue] (0.1,0.5) circle (1.2pt);

% real line
\draw[realline, targetgreen] (5.2,-1.3) -- (5.2,1.3);
\foreach \y in {-1, 0, 1} {
  \draw[targetgreen] (5.05,\y) -- (5.35,\y);
}
\node[targetgreen] at (5.2,1.7) {$(\R,\mathcal{B}(\R))$};
\node[targetgreen] at (5.2,-1.7) {real line};
\filldraw[targetgreen] (5.2,0.25) circle (1.4pt) node[right, font=\tiny]{$X(\omega)$};

% map X
\draw[arr, quotientorange] (-0.55,0.15) to[bend left=18] (5.1,0.27);
\node[quotientorange] at (2.4,1.15) {$X$ (measurable)};

% pushforward law
\draw[arr, structurepurple] (1.6,-0.9) to[bend right=22] (4.6,-0.9);
\node[structurepurple] at (3.0,-1.05) {$\P_X=\P\circ X^{-1}$};
\end{tikzpicture}
\caption{The structure of a probability space. A random variable $X$ is a measurable map
from $(\Omega,\mathcal{F},\P)$ to $(\R, \mathcal{B}(\R))$. Measurability ensures that
the preimage of every Borel set is a valid event in $\mathcal{F}$. The pushforward
$\P_X = \P \circ X^{-1}$ is the distribution of $X$, which determines all probabilities
and expectations on $\R$.}
\label{ch36:fig-anatomy}
\end{figure}

\begin{recoverybox}
Those three definitions --- $\sigma$-algebra, measure, and measurable function --- are
the entire grammar of a probability space. None of the calculations from Chapters~1--35
change; they simply gain a single, contradiction-free home.
\end{recoverybox}

\section{Summary of measure-theoretic foundations}
\label{ch36:payoff}

Measure theory provides four primary structural advantages for probability.
\emph{First}, consistency: it establishes a rigorous framework for continuous sample
spaces where unrestricted probability functions lead to contradictions (Section~\ref{ch36:crack}).
\emph{Second}, unification: discrete sums, continuous integrals, and mixed distributions
are unified under a single definition of integration against a measure (Section~\ref{ch36:unify}).
\emph{Third}, limits: countable additivity ensures continuity along monotone sequences of events
(Proposition~\ref{ch36:props}), justifying limit operations in asymptotic theorems and
infinite-dimensional stochastic processes.
\emph{Fourth}, information: sub-$\sigma$-algebras provide a precise formalization of
conditioning, martingales, and evolving information structures.
Next, in Chapter~37, we put $(\Omega, \mathcal{F}, \P)$ to work by constructing the Lebesgue integral.

\section{Exercises}
\label{ch36:exercises}

\begin{exercise}[Warm-up mental checks]
\label{ch36:ex-warmup}
Answer these quick checks before reaching for paper:
\begin{enumerate}[label=(\alph*)]
  \item On any sample space $\Omega$, what are the coarsest and finest possible $\sigma$-algebras?
  \item Why does a $\sigma$-algebra require closure under \emph{countable} unions rather than finite or arbitrary unions?
  \item If an event $N$ has $\P(N) = 0$, is it logically impossible for an outcome in $N$ to occur in a continuous experiment?
  \item For a real-valued random variable $X$, what is the pushforward measure $\P_X$ evaluated on the half-line $(-\infty, x]$ called in elementary probability?
\end{enumerate}
\end{exercise}

\begin{exercise}[Warm-up with the axioms]
\label{ch36:ex-axioms}
Let $\Omega = \set{1,2,3,4}$. Decide which of the following are
$\sigma$-algebras on $\Omega$, and for those that are not, name the failing axiom.
(a) $\set{\emptyset, \Omega}$.
(b) $\set{\emptyset, \set{1}, \set{2,3,4}, \Omega}$.
(c) $\set{\emptyset, \set{1}, \set{2}, \Omega}$.
(d) The collection of all subsets with an even number of elements.
\end{exercise}

\begin{exercise}[Generated $\sigma$-algebra]
\label{ch36:ex-generated}
On $\Omega = \set{1,2,3,4,5,6}$ (a die), write down the smallest
$\sigma$-algebra containing the events $A = \set{1,2}$ and $B = \set{2,3}$. List
all its members. How many are there? (Hint: it is generated by the ``atoms'' you
can no longer split using $A$, $B$, and complements.)
\end{exercise}

\begin{exercise}[Intersections preserve structure]
\label{ch36:ex-intersection}
Prove that if $\mathcal{F}_1$ and $\mathcal{F}_2$ are $\sigma$-algebras on the same
$\Omega$, then $\mathcal{F}_1 \cap \mathcal{F}_2$ is a $\sigma$-algebra. Then give
an example showing $\mathcal{F}_1 \cup \mathcal{F}_2$ need \emph{not} be. (This is
why ``generated by'' is defined via intersection, not union.)
\end{exercise}

\begin{exercise}[Countable sets are null]
\label{ch36:ex-null}
Using Lebesgue measure $\lambda$ on $\R$, prove $\lambda(\Q) = 0$ by covering the
$n$th rational with an interval of length $\eps / 2^n$ and using countable
subadditivity. Then explain in one sentence why this does \emph{not} let you
conclude $\lambda([0,1]) = 0$, even though $[0,1]$ is also a union of single
points.
\end{exercise}

\begin{exercise}[Continuity from above]
\label{ch36:ex-cont-above}
State and prove the ``continuity from above'' analogue of
Proposition~\ref{ch36:props}(4): if $A_1 \supseteq A_2 \supseteq \cdots$ and
$\P(A_1) < \infty$, then $\P\!\left(\bigcap_n A_n\right) = \lim_n \P(A_n)$.
(Hint: apply continuity from below to the complements $A_n^c$.) Give an example on
$\R$ with Lebesgue measure where dropping the finiteness assumption makes the
conclusion fail.
\end{exercise}

\begin{exercise}[Measurability via half-lines]
\label{ch36:ex-measurable}
Let $X : \Omega \to \R$ satisfy $\set{X \le x} \in \mathcal{F}$ for every rational
$x$. Prove that $\set{X \le x} \in \mathcal{F}$ for every \emph{real} $x$, hence
$X$ is a random variable. (Hint: write $\set{X \le x}$ as a countable
intersection over rationals decreasing to $x$, and use closure.)
\end{exercise}

\begin{exercise}[Pushforward is a measure]
\label{ch36:ex-pushforward}
Let $X$ be a random variable on $(\Omega, \mathcal{F}, \P)$ and define
$\P_X(B) = \P(X^{-1}(B))$ for Borel $B$. Verify directly that $\P_X$ satisfies the
two axioms of a probability measure (Definition~\ref{ch36:measure-def} plus
normalization). Which property of preimages do you need for countable additivity?
\end{exercise}

\begin{exercise}[A mixed distribution]
\label{ch36:ex-mixed}
For the crosswalk waiting time $T$ of Example~\ref{ch36:mixed} with law
$\P_T = \tfrac12\delta_0 + \tfrac12\lambda|_{(0,1)}$, compute (a) $\P(T = 0)$;
(b) $\P(T > 1/2)$; (c) the CDF $F_T(t)$ for all $t$; and (d) $\E[T]$, treating the
expectation as $\tfrac12 \cdot 0 + \tfrac12 \int_0^1 t\,\dd t$. Sketch $F_T$ and
identify the jump.
\end{exercise}

\begin{exercise}[Almost surely versus surely]
\label{ch36:ex-as}
Let $U \sim \Unif(0,1)$. Give an event that is (a) impossible (empty);
(b) possible but of probability $0$; (c) certain to hold (all of $\Omega$);
(d) of probability $1$ but \emph{not} all of $\Omega$. Explain why (b) and (d) are
two sides of the same phenomenon, and why no analogue of (b) or (d) exists when
$\Omega$ is finite.
\end{exercise}

\begin{exercise}[Why countable, not arbitrary --- a proof]
\label{ch36:ex-why-countable}
Suppose, for contradiction, that we strengthened
Definition~\ref{ch36:sigma-def} to require closure under \emph{arbitrary} unions.
Show that the only such ``$\sigma$-algebra'' on $\R$ containing every singleton
$\set{x}$ is the full power set $2^{\R}$. Combine this with
Theorem~\ref{ch36:vitali} to explain in a sentence why arbitrary-union closure is
incompatible with a translation-invariant uniform probability --- i.e.\ why the
word ``countable'' in the definition is not negotiable.
\end{exercise}

\begin{exercise}[Conceptual: information as a $\sigma$-algebra]
\label{ch36:ex-info}
Two fair coins are flipped, so $\Omega = \set{HH, HT, TH, TT}$ with the power set and
uniform $\P$. Let $\mathcal{G}$ be the smallest $\sigma$-algebra making the
random variable ``number of heads'' measurable. (a) List the members of
$\mathcal{G}$. (b) Is the indicator for ``did the \emph{first} coin land heads?''
measurable with respect to $\mathcal{G}$? Explain what this says about which
questions ``knowing only the number of heads'' lets you answer.
\end{exercise}

\begin{exercise}[Proof: agreement a.s.\ implies equal law]
\label{ch36:ex-as-law}
Let $X, Y$ be random variables on the same space with $\P(X = Y) = 1$. Prove that
$\P_X = \P_Y$, i.e.\ $\P(X \in B) = \P(Y \in B)$ for every Borel set $B$. (Hint:
the events $\set{X \in B}$ and $\set{Y \in B}$ differ only within the null set
$\set{X \ne Y}$; bound the symmetric difference.) Conclude that $X$ and $Y$ have
the same CDF, expectation, and variance --- the precise sense in which measure
theory ``cannot tell them apart.''
\end{exercise}
